\chapter{L’économie des sépultures}\label{ch:sepultures}

\section{Acheter une tombe}

L’ensemble B contient sept actes d’achat ou de réservation d’emplacement sans commémoration (\nnums{\st{acteslist}}) et sept textes mixtes, qui associent une acquisition et la commémoration d’un défunt (\nnums{\st{mixteslist}} ; tableau~\ref{tab:achats}). Ces quatorze textes sont tous postérieurs à 369. Ils se concentrent dans les deux dernières décennies : dix d’entre eux datent des années 385–403.

% généré par scripts/04_generated.py
\begin{longtable}{@{}r r l p{0.68\linewidth}@{}}\caption{Actes d’achat ou de réservation d’emplacement et textes mixtes (ensemble B ; début du texte transmis).}\label{tab:achats}\\\toprule N° & Année & Genre & Texte (début)\\\midrule\endfirsthead\toprule N° & Année & Genre & Texte (début)\\\midrule\endhead\bottomrule\endfoot
137 & 370 & mixte & {\small\lat{[- - -] et Artemia / [locum se vivi si]bi cum suis / [comparaverunt] d(e)p(osita) Artemia XIIII Kal(endas) Sept(embres) / [Valentin(iano) III et Val]e(n)t(is) III Augg(ustis) conss(ulibus)}} \\
141 & 370 & acte & {\small\lat{Concordius se / vivo locum bisomum / emit c(onsule) Valentiniano III}} \\
209 & 378 & acte & {\small\lat{Hic locus / est Victo / rini quem / emit a Mar / tinian⸢o⸣ / fossore / Valente VI / et Valentin(iano !)}} \\
217 & 380 & acte & {\small\lat{Gratiano V et Th(e)odosio [conss(ulibus)] / ego Fl(avius) Victor me vi⸢v⸣o s[cripsi] / co⸢m⸣parasse cum Cons[- - -] / [coniuge me]a locum / [bisomum a La]ur() fossorem}} \\
281 & 385 & acte & {\small\lat{Locus Rufin[ae quae] / se ⸢v⸣i⸢v⸣a sibi c⸢o⸣[mpa] / ra⸢v⸣it bisomu(m) [d(omino) n(ostro) Arcad] / io et Baudone v(iro) [c(larissimo) cons(ulibus)]}} \\
340 & 391 & acte & {\small\lat{Fl(avio) Tatiano et Quinto / Aur(elio) S⸢y⸣mmaco v⸢i⸣ris / cl(a)rissimis ego Zita / locum quadri\{c\} / somu(m) in b(a)silic(a) / al⸢b⸣a emi[t]}} \\
346 & 392 & mixte & {\small\lat{[C]⸢o⸣mparaverun[t sibi locum] / se vivi Iustinus et Feli[citas] / dep(osita) Pitzinnina in pace R[- - -] / qu(a)e vixit ann(os) XXVI et [- - -] / Arcadio Aug(usto) II et Rufin[o cons(ulibus…}} \\
387 & 396 & mixte & {\small\lat{[- - -] deposita X[- - - d(omino) n(ostro) Honorio Aug(usto) III et] / d(omino) n(ostro) Fl(avio) Ar(ca)dio Aug(usto) IIII conss(ulibus) M[- - -] / Eros et Leo et Setema locu(m) [emerunt ad…}} \\
400 & 397 & mixte & {\small\lat{Dulcis et innoce(n)s hic dormit Severianus in somn\{i\}o pacis / qui vixit ann⸢o⸣s p(lus) m(inus) L cuius spiritus in luce domini suscep / tus est die VIIIIX(!) K(a)l(endas) Mart(ias) dd(omini…}} \\
406 & 397 & mixte & {\small\lat{Sussanna compara / vit sibi memoriam quie / vit die VII Kal(endas) Augusta⸢s⸣ / conss(ulibus) Caesario et Attico}} \\
411 & 397 & mixte & {\small\lat{Deposita puella nomine / Ursa in pace annorum / viginti duo conss(ulatu) Fl(avi) Caesari et Maximi / vv(irorum) cc(larissimorum) emptu(m) locu(m) (t)r⸢e⸣s et s(olidis) / XI}} \\
417 & 398 & acte & {\small\lat{Dom[inis nos]tris Ho(no)rio quater et Eut⸢y⸣c(h)iano consulibus / loc[um Felici]tati(s) et Crescentiaes(!) [- - -] em[erunt s]e ⸢v⸣i⸢v⸣i sibi et suis}} \\
424 & 398 & acte & {\small\lat{Locus Zosimi et Concordi(a)es(!) / co⸢m⸣paratus a Preiecticio / fossore / consulatus domi(no) nostro Ho / norio quater et Eut⸢y⸣c(h)iano consulib(us)}} \\
478 & 403 & mixte & {\small\lat{Rimorido c(larissimo) c(onsule) mense April(i) XVIII Kal(endas) / Maias Barbarus defunctus lo / cu(s) bisomu(s) em(p)tu(s) ab Urs⸢o⸣ fossore}} \\
\end{longtable}


Les emplacements sont acquis par des particuliers, souvent de leur vivant (\emph{se vivo}, \emph{se viva}, \emph{se vivi} : \nnums{137}, 141, 217, 281, 346 et 417), et pour un couple ou une famille (\emph{bisomus} aux \nnums{141}, 217, 281 et 478 ; \emph{quadrisomus} au \nnum{340} ; \emph{sibi et suis} au \nnum{417} ; \emph{sibi cum suis} au \nnum{137}). Le verbe est \emph{emere} ou \emph{comparare}. Flavius Victor rédige l’acte à la première personne : \emph{ego Fl(avius) Victor me vivo s[cripsi] comparasse cum Cons[- - -] [coniuge me]a locum [bisomum]} (\nnum{217}, 380). Zita fait de même, \emph{ego Zita locum quadrisomum in basilica alba emit} (\nnum{340}, 391). L’expression \emph{basilica alba} (EDB : \emph{alva}) désigne vraisemblablement un édifice du cimetière de Domitilla, auquel l’EDB rattache le texte ; son identification n’est pas proposée ici.

\section{Les vendeurs}

Quatre textes nomment le vendeur, toujours un \emph{fossor} : Martinianus (\nnum{209}, 378), Laur[- - -] (\nnum{217}, 380 ; le nom est abrégé et en partie restitué), Preiecticius (\nnum{424}, 398) et Ursus (\nnum{478}, 403). Ces textes confirment, à l’échelle de la Rome de la fin du \textsc{iv}\textsuperscript{e} siècle, le rôle des fossoyeurs dans la vente des emplacements.

Deux textes ont été interprétés comme des achats auprès du clergé ; aucun ne le prouve.

\emph{Le \nnum{387} (396).} Le texte transmis par l’EDCS est : \emph{Eros et Leo et Setema locu(m) [emerunt ad limina Petri et] / Pauli sub episcopatu Si[rici a presbb(yteris) Urso] / et Proclino}. Sont lus : les noms des acquéreurs, \emph{locum}, \emph{Pauli sub episcopatu Si-} et \emph{et Proclino}. Le verbe, la mention de Pierre, les \emph{presbyteri} et le nom d’Ursus sont restitués. Le texte n’est pas apparié à l’EDB. La datation \emph{sub episcopatu Si[rici]} est vraisemblable, car l’année consulaire est 396 et Sirice est évêque de Rome de 384 à 399 ; la localisation près de la tombe de Paul et la qualité de prêtres des vendeurs dépendent de restitutions.

\emph{Le \nnum{400} (397).} Viricunda a acheté l’emplacement de son époux Severianus, \emph{quem locum emit Viricunda uxor Pascasio ep(is)c(opo)}. L’EDB ne développe pas l’abréviation (\emph{EPC}) et reprend la note de Diehl : aucun évêque de Rome ne porte le nom de Pascasius, et le quantième transmis, \emph{VIIIIX K(a)l(endas) Mart(ias)}, serait à corriger. Le datif \emph{Pascasio} peut désigner le vendeur, le bénéficiaire ou l’autorité sous laquelle la transaction a eu lieu. La fonction de ce Pascasius reste indéterminée.

\section{Le prix}

Le \nnum{411} (397) est le seul texte qui indique un prix. L’EDCS lit \emph{emptu(m) locu(m) (t)r⸢e⸣s et s(olidis) XI}, l’EDB \emph{emptu locu ⟨t⟩ris et s(olidis) (undecim)}. Les deux éditions suppléent une lettre initiale omise et interprètent \emph{s.} comme \emph{solidis}. Si la lecture est exacte, le prix est de onze sous d’or, ou de trois et onze sous selon la construction de \emph{tres} ; le texte ne permet pas d’en décider. Le collège consulaire transmis par ce texte (\emph{Fl(avi) Caesari et Maximi}) ne correspond pas aux fastes de 397, et la notice reste sous réserve.

\section{Les textes mixtes}

Dans les textes mixtes, l’acquisition et la commémoration se succèdent sur la même plaque. Au \nnum{406}, Susanna « s’est acheté une \emph{memoria} » et « s’est reposée » le 26 juillet 397 : la date se rapporte au décès, non à l’achat. Au \nnum{346}, Iustinus et Felicitas ont acheté de leur vivant l’emplacement où est déposée Pitzinnina. Au \nnum{478}, le défunt Barbarus repose dans un \emph{locus bisomus} acheté au fossoyeur Ursus. Ces textes montrent que l’inscription de la propriété de la tombe et celle du défunt relevaient d’une même pratique ; la distinction entre « acte » et « épitaphe », nécessaire au calcul des dénominateurs, est en partie une construction moderne (chapitre~\ref{ch:corpus}).

\section{\emph{Ad sanctos}}

Deux textes témoignent de la recherche d’une sépulture privilégiée. Le \nnum{232} (382) dit que la défunte, pour ses mérites, a reçu une sépulture \emph{[- - -] limina sanctorum}, « ce que [beaucoup] désirent et que peu obtiennent » (\emph{[quod multi cupiu]n(t) et rari accipiun(t)}) : \emph{limina sanctorum} et la seconde moitié de la formule sont lus ; le premier mot de l’expression et l’antithèse \emph{multi / rari} sont restitués. Le texte ne dit pas quels saints, ni où. Il est d’origine incertaine. Au \nnum{387}, seul \emph{Pauli} est lu de l’expression restituée \emph{[ad limina Petri et] Pauli}.

Ces textes sont compatibles avec la thèse d’Éric Rebillard sur la responsabilité des familles dans la sépulture et sur le caractère tardif et limité de l’encadrement ecclésiastique des cimetières\footnote{É.~Rebillard, \emph{Religion et sépulture. L’Église, les vivants et les morts dans l’Antiquité tardive}, Paris, 2003 (Civilisations et sociétés, 115) ; trad. angl. \emph{The Care of the Dead in Late Antiquity}, Ithaca, 2009. Notice vérifiée ; aucun renvoi de page.}, sans permettre de la vérifier dans sa généralité. L’intervention du clergé dans la vente des tombes n’est attestée, dans ce corpus, que par des restitutions.
